\documentclass[12pt,a4paper]{report}
\usepackage[utf8]{inputenc}
\usepackage[T1]{fontenc}
\usepackage[french]{babel}
\usepackage{geometry}
\usepackage{graphicx}
\usepackage{listings}
\usepackage{xcolor}
\usepackage{hyperref}
\usepackage{longtable}
\usepackage{array}
\usepackage{booktabs}
\usepackage{fancyhdr}
\usepackage{titlesec}
\usepackage{tocloft}
\usepackage{amsmath}
\usepackage{enumitem}
\usepackage{float}
\usepackage{caption}
\usepackage{tikz}
\usetikzlibrary{shapes,arrows,arrows.meta,positioning,fit,calc}

\geometry{margin=2.5cm}

% Couleurs pour listings
\definecolor{codebg}{RGB}{245,245,245}
\definecolor{codeframe}{RGB}{200,200,200}
\definecolor{keyword}{RGB}{0,0,150}
\definecolor{comment}{RGB}{0,128,0}
\definecolor{string}{RGB}{163,21,21}

\lstset{
  basicstyle=\ttfamily\small,
  backgroundcolor=\color{codebg},
  frame=single,
  rulecolor=\color{codeframe},
  breaklines=true,
  breakatwhitespace=true,
  tabsize=2,
  showstringspaces=false,
  keywordstyle=\color{keyword}\bfseries,
  commentstyle=\color{comment}\itshape,
  stringstyle=\color{string},
  numbers=left,
  numberstyle=\tiny\color{gray},
  numbersep=5pt,
  captionpos=b
}

\lstdefinelanguage{DSL}{
  keywords={LoadBalancerSystem, cluster, server, loadbalancer, healthCheck, scalingRule, listener, alert, sessionPersistence},
  keywordstyle=\color{keyword}\bfseries,
  morecomment=[l]{//},
  morestring=[b]",
  sensitive=true
}

\lstdefinelanguage{OCL}{
  keywords={context, inv, self, implies, forAll, exists, select, collect, notEmpty, isUnique, size, and, or, not},
  keywordstyle=\color{keyword}\bfseries,
  morecomment=[l]{--},
  sensitive=true
}

\lstdefinelanguage{ATL}{
  keywords={module, create, from, rule, to, lazy, helper, context, def, if, then, else, endif, let, in},
  keywordstyle=\color{keyword}\bfseries,
  morecomment=[l]{--},
  morestring=[b]',
  sensitive=true
}

\lstdefinelanguage{Xtext}{
  keywords={grammar, with, generate, enum, terminal, returns, import},
  keywordstyle=\color{keyword}\bfseries,
  morecomment=[l]{//},
  morestring=[b]",
  sensitive=true
}

% En-têtes
\pagestyle{fancy}
\fancyhf{}
\fancyhead[L]{\small INF5039 -- Projet IDM}
\fancyhead[R]{\small Répartition de charge par DSL}
\fancyfoot[C]{\thepage}

\hypersetup{
  colorlinks=true,
  linkcolor=black,
  urlcolor=blue,
  citecolor=black
}

\begin{document}

% ============================================================
% PAGE DE GARDE
% ============================================================
\begin{titlepage}
\centering
\vspace*{1cm}

{\Large\bfseries UNIVERSITÉ DE YAOUNDÉ 1}\\[0.3cm]
{\large Faculté des Sciences}\\[0.2cm]
{\large Département d'Informatique}\\[0.2cm]
{\large Master 2 -- INF5039 : Ingénierie Dirigée par les Modèles}\\[2cm]

\rule{\textwidth}{1pt}\\[0.5cm]
{\LARGE\bfseries PROJET}\\[0.5cm]
{\Large\bfseries Modélisation et génération d'un mécanisme\\de répartition de charge par DSL\\pour un système distribué}\\[0.5cm]
\rule{\textwidth}{1pt}\\[2.5cm]

{\large\bfseries Réalisé par :}\\[0.3cm]
{\large ZOGO ABOUMA Zozime Achaire -- 18N2824}\\[0.2cm]
{\large NKOLO Stéphane Roylex -- 21T2588}\\[2cm]

{\large\bfseries Encadreur :}\\[0.3cm]
{\large Dr MONTHE Valérie}\\[0.2cm]
{\large Université de Yaoundé 1}\\[2.5cm]

{\large\bfseries Année académique 2025 -- 2026}

\end{titlepage}

% ============================================================
% TABLE DES MATIERES
% ============================================================
\tableofcontents
\newpage

% ============================================================
% INTRODUCTION
% ============================================================
\chapter{Introduction}

Dans le contexte actuel des systèmes informatiques, la montée en charge des applications distribuées impose des mécanismes efficaces de répartition de charge (\textit{load balancing}). Ces mécanismes permettent de distribuer équitablement les requêtes entre plusieurs serveurs afin d'optimiser les performances, d'assurer la haute disponibilité et de prévenir la surcharge de nœuds individuels.

L'Ingénierie Dirigée par les Modèles (IDM/MDE) offre une approche systématique pour modéliser, transformer et générer des artefacts logiciels à partir de modèles abstraits. En combinant IDM et la définition d'un langage dédié (DSL -- \textit{Domain Specific Language}), il devient possible de spécifier de manière déclarative les politiques de répartition de charge et de générer automatiquement le code d'infrastructure correspondant.

Ce projet a pour objectif de concevoir un DSL permettant de modéliser un mécanisme de répartition de charge pour un système distribué, puis de mettre en œuvre une chaîne de transformations MDA (\textit{Model Driven Architecture}) allant du modèle abstrait jusqu'au code exécutable.

% ============================================================
% CHAPITRE 1 : DESCRIPTION DU DOMAINE
% ============================================================
\chapter{Description du domaine traité}

\section{Qu'est-ce que la répartition de charge ?}

La répartition de charge (\textit{load balancing}) est une technique utilisée dans les systèmes distribués pour distribuer le trafic réseau ou les tâches de calcul entre plusieurs serveurs ou ressources. L'objectif est de :

\begin{itemize}
  \item Maximiser le débit (\textit{throughput}) du système global
  \item Minimiser le temps de réponse pour chaque requête
  \item Éviter la surcharge d'un seul nœud
  \item Assurer la tolérance aux pannes (\textit{failover})
  \item Permettre la scalabilité horizontale
\end{itemize}

\section{Composants d'un système de répartition de charge}

Un système de répartition de charge typique comprend les éléments suivants :

\begin{itemize}
  \item \textbf{Load Balancer} (répartiteur) : point d'entrée qui reçoit les requêtes et les redirige
  \item \textbf{Backend Servers} (serveurs backend) : ensemble de serveurs qui traitent les requêtes
  \item \textbf{Health Check} (vérification de santé) : mécanisme pour vérifier la disponibilité des serveurs
  \item \textbf{Algorithme de répartition} : stratégie utilisée pour choisir le serveur cible
  \item \textbf{Session Persistence} (affinité de session) : mécanisme pour diriger un même client vers le même serveur
  \item \textbf{Cluster} : regroupement logique de serveurs backend
  \item \textbf{Listeners} : points d'écoute sur des ports/protocoles spécifiques
\end{itemize}

\section{Algorithmes de répartition courants}

\begin{table}[H]
\centering
\caption{Algorithmes de répartition de charge}
\begin{tabular}{|p{4cm}|p{5.5cm}|p{4.5cm}|}
\hline
\textbf{Algorithme} & \textbf{Description} & \textbf{Cas d'usage} \\
\hline
Round Robin & Distribution séquentielle cyclique entre les serveurs & Serveurs homogènes \\
\hline
Weighted Round Robin & Round Robin avec poids assignés & Serveurs hétérogènes \\
\hline
Least Connections & Envoie au serveur ayant le moins de connexions actives & Requêtes de durée variable \\
\hline
IP Hash & Hash de l'IP source pour déterminer le serveur & Affinité de session \\
\hline
Random & Choix aléatoire & Cas simples, tests \\
\hline
Weighted Least Conn. & Least Connections pondéré par la capacité & Serveurs hétérogènes, charge variable \\
\hline
\end{tabular}
\end{table}

\section{Contexte et problématique}

Traditionnellement, la configuration d'un load balancer se fait manuellement via des fichiers de configuration (nginx.conf, haproxy.cfg, etc.) qui sont :

\begin{itemize}
  \item Spécifiques à une technologie (\textit{vendor lock-in})
  \item Sujets à des erreurs de syntaxe
  \item Difficiles à valider avant déploiement
  \item Non portables d'une plateforme à l'autre
\end{itemize}

Notre approche IDM propose de résoudre ces problèmes en définissant un DSL indépendant de toute technologie qui capture l'essence de la répartition de charge, et en utilisant des transformations de modèles pour générer automatiquement les configurations ou le code cible.

% ============================================================
% CHAPITRE 2 : CONCEPTS
% ============================================================
\chapter{Tableau des concepts identifiés}

Le tableau suivant recense les concepts clés du domaine de la répartition de charge qui seront modélisés dans notre métamodèle :

\begin{longtable}{|c|p{3.5cm}|p{5cm}|p{4.5cm}|}
\hline
\textbf{N°} & \textbf{Concept} & \textbf{Description} & \textbf{Attributs} \\
\hline
\endhead
1 & LoadBalancerSystem & Système global, élément racine & name, description, version \\
\hline
2 & Cluster & Groupe logique de serveurs backend & name, maxConnections \\
\hline
3 & Server & Serveur backend individuel & name, host, port, weight, maxConnections, enabled \\
\hline
4 & LoadBalancer & Composant répartiteur & name, algorithm, stickySession \\
\hline
5 & Algorithm & Énumération des algorithmes & ROUND\_ROBIN, WEIGHTED\_RR, LEAST\_CONN, IP\_HASH, RANDOM \\
\hline
6 & HealthCheck & Vérification de disponibilité & protocol, path, interval, timeout, thresholds \\
\hline
7 & Listener & Point d'écoute réseau & name, protocol, port, targetCluster \\
\hline
8 & Protocol & Protocoles supportés & HTTP, HTTPS, TCP, UDP \\
\hline
9 & SessionPersistence & Affinité de session & type, cookieName, ttl \\
\hline
10 & Alert & Règle d'alerte & name, metric, threshold, action \\
\hline
11 & ScalingRule & Mise à l'échelle auto & name, metric, scaleUp/Down, min/maxInstances \\
\hline
12 & Metric & Métriques surveillées & CPU, MEMORY, REQUESTS, RESPONSE\_TIME, CONNECTIONS \\
\hline
\end{longtable}

% ============================================================
% CHAPITRE 3 : REGLES DE GESTION
% ============================================================
\chapter{Règles de gestion et exigences}

Les règles suivantes doivent être respectées dans tout modèle conforme à notre métamodèle :

\begin{longtable}{|c|p{4cm}|p{8.5cm}|}
\hline
\textbf{Code} & \textbf{Règle} & \textbf{Description} \\
\hline
\endhead
RG01 & Unicité des noms de serveurs & Tous les serveurs d'un même cluster doivent avoir des noms uniques \\
\hline
RG02 & Port valide & Le port doit être compris entre 1 et 65535 \\
\hline
RG03 & Poids positif & Le poids d'un serveur doit être $\geq 1$ \\
\hline
RG04 & Cluster non vide & Un cluster doit contenir au moins un serveur \\
\hline
RG05 & Health check cohérent & L'intervalle doit être supérieur au timeout \\
\hline
RG06 & Seuils de scaling cohérents & Le seuil de scale-up $>$ seuil de scale-down \\
\hline
RG07 & Instances min/max & minInstances $\leq$ maxInstances \\
\hline
RG08 & Listener unique par port & Deux listeners ne peuvent pas écouter sur le même port \\
\hline
RG09 & Weighted requiert des poids & Si algorithme = Weighted RR, tous les serveurs doivent avoir un poids $\geq 1$ \\
\hline
RG10 & Cookie name requis & Si persistence = COOKIE, cookieName doit être renseigné \\
\hline
RG11 & Au moins un listener & Le système doit avoir au moins un listener \\
\hline
RG12 & Serveur actif requis & Un cluster doit avoir au moins un serveur enabled = true \\
\hline
\end{longtable}

% ============================================================
% CHAPITRE 4 : DECOMPOSITION EN MODULES
% ============================================================
\chapter{Décomposition en modules}

\section{Architecture des packages}

Notre système est organisé en modules pour une meilleure séparation des préoccupations :

\begin{lstlisting}[language={},caption={Structure en packages du projet}]
loadbalancer.dsl/
  metamodel/                    # Metamodele
    core/                       # Concepts fondamentaux
      LoadBalancerSystem, Cluster, Server
    balancing/                  # Repartition
      LoadBalancer, Algorithm, SessionPersistence
    networking/                 # Reseau
      Listener, Protocol
    monitoring/                 # Surveillance
      HealthCheck, Alert, Metric
    scaling/                    # Mise a l'echelle
      ScalingRule
  constraints/                  # Contraintes OCL
  concrete.syntax/              # Syntaxe concrete (Xtext)
  transformations/
    m2m/                        # Modele-a-modele (ATL)
    m2t/                        # Modele-a-texte (Acceleo)
  examples/                     # Etudes de cas
\end{lstlisting}

\section{Description des modules}

\begin{table}[H]
\centering
\caption{Description des modules du système}
\begin{tabular}{|p{3.5cm}|p{5cm}|p{5cm}|}
\hline
\textbf{Module} & \textbf{Rôle} & \textbf{Contenu} \\
\hline
metamodel.core & Concepts fondamentaux & LoadBalancerSystem, Cluster, Server \\
\hline
metamodel.balancing & Logique de répartition & LoadBalancer, Algorithm, SessionPersistence \\
\hline
metamodel.networking & Aspects réseau & Listener, Protocol \\
\hline
metamodel.monitoring & Surveillance et alertes & HealthCheck, Alert, Metric \\
\hline
metamodel.scaling & Mise à l'échelle & ScalingRule \\
\hline
constraints & Contraintes OCL & Fichiers .ocl \\
\hline
concrete.syntax & Grammaire Xtext & Fichier .xtext \\
\hline
transformations.m2m & Transfo M2M & Fichiers .atl \\
\hline
transformations.m2t & Transfo M2T & Fichiers .mtl \\
\hline
\end{tabular}
\end{table}

% ============================================================
% CHAPITRE 5 : METAMODELE
% ============================================================
\chapter{Schéma du métamodèle}

\section{Diagramme global du métamodèle}

\begin{figure}[H]
\centering
\begin{tikzpicture}[
  class/.style={rectangle, draw, fill=white, minimum width=4cm, minimum height=1.5cm, align=center, font=\small},
  enum/.style={rectangle, draw, fill=gray!10, minimum width=3cm, minimum height=1.2cm, align=center, font=\small},
  arrow/.style={-{Stealth}, thick},
  comp/.style={-{Diamond[open]}, thick}
]

% Classes principales
\node[class] (system) at (0,0) {\textbf{LoadBalancerSystem}\\name, description, version};
\node[class] (cluster) at (-4,-3) {\textbf{Cluster}\\name, maxConnections};
\node[class] (listener) at (4,-3) {\textbf{Listener}\\name, protocol, port};
\node[class] (server) at (-6,-6) {\textbf{Server}\\name, host, port, weight, enabled};
\node[class] (lb) at (-1,-6) {\textbf{LoadBalancer}\\algorithm, stickySession};
\node[class] (hc) at (-6,-9) {\textbf{HealthCheck}\\protocol, path, interval, timeout};
\node[class] (scaling) at (-1,-9) {\textbf{ScalingRule}\\metric, thresholds, min/max};
\node[class] (alert) at (4,-6) {\textbf{Alert}\\name, metric, threshold, action};
\node[class] (sp) at (4,-9) {\textbf{SessionPersistence}\\type, cookieName, ttl};

% Enums
\node[enum] (algo) at (-1,-11.5) {\textbf{<<enum>> Algorithm}};
\node[enum] (proto) at (4,-11.5) {\textbf{<<enum>> Protocol}};
\node[enum] (metric) at (-6,-11.5) {\textbf{<<enum>> Metric}};

% Relations
\draw[comp] (system) -- node[left, font=\scriptsize] {1..*} (cluster);
\draw[comp] (system) -- node[right, font=\scriptsize] {1..*} (listener);
\draw[comp] (system) -| node[right, font=\scriptsize] {0..*} (alert);
\draw[comp] (cluster) -- node[left, font=\scriptsize] {1..*} (server);
\draw[comp] (cluster) -- node[right, font=\scriptsize] {1} (lb);
\draw[comp] (cluster) -- node[left, font=\scriptsize] {0..1} (hc);
\draw[comp] (cluster) -- node[right, font=\scriptsize] {0..1} (scaling);
\draw[comp] (lb) -- node[right, font=\scriptsize] {0..1} (sp);
\draw[arrow, dashed] (listener) -- node[above, font=\scriptsize] {targetCluster} (cluster);

% Enum links
\draw[arrow, dotted] (lb) -- (algo);
\draw[arrow, dotted] (listener) -- (proto);
\draw[arrow, dotted] (scaling) -- (metric);

\end{tikzpicture}
\caption{Diagramme de classes du métamodèle global}
\end{figure}

\section{Module Core -- Détail}

\begin{lstlisting}[language=XML, caption={Extrait du métamodèle Ecore (loadbalancer.ecore)}]
<?xml version="1.0" encoding="UTF-8"?>
<ecore:EPackage xmi:version="2.0"
    xmlns:xmi="http://www.omg.org/XMI"
    xmlns:xsi="http://www.w3.org/2001/XMLSchema-instance"
    xmlns:ecore="http://www.eclipse.org/emf/2002/Ecore"
    name="loadbalancer"
    nsURI="http://www.example.org/loadbalancer"
    nsPrefix="lb">

  <eClassifiers xsi:type="ecore:EClass" name="LoadBalancerSystem">
    <eStructuralFeatures xsi:type="ecore:EAttribute" name="name"
        eType="ecore:EDataType http://www.eclipse.org/emf/2002/Ecore#//EString"/>
    <eStructuralFeatures xsi:type="ecore:EAttribute" name="description"
        eType="ecore:EDataType http://www.eclipse.org/emf/2002/Ecore#//EString"/>
    <eStructuralFeatures xsi:type="ecore:EAttribute" name="version"
        eType="ecore:EDataType http://www.eclipse.org/emf/2002/Ecore#//EString"/>
    <eStructuralFeatures xsi:type="ecore:EReference" name="clusters"
        upperBound="-1" eType="#//Cluster" containment="true"/>
    <eStructuralFeatures xsi:type="ecore:EReference" name="listeners"
        upperBound="-1" eType="#//Listener" containment="true"/>
    <eStructuralFeatures xsi:type="ecore:EReference" name="alerts"
        upperBound="-1" eType="#//Alert" containment="true"/>
  </eClassifiers>

  <eClassifiers xsi:type="ecore:EClass" name="Cluster">
    <eStructuralFeatures xsi:type="ecore:EAttribute" name="name"
        eType="ecore:EDataType http://www.eclipse.org/emf/2002/Ecore#//EString"/>
    <eStructuralFeatures xsi:type="ecore:EReference" name="servers"
        upperBound="-1" eType="#//Server" containment="true"/>
    <eStructuralFeatures xsi:type="ecore:EReference" name="loadBalancer"
        eType="#//LoadBalancer" containment="true"/>
    <eStructuralFeatures xsi:type="ecore:EReference" name="healthCheck"
        eType="#//HealthCheck" containment="true"/>
    <eStructuralFeatures xsi:type="ecore:EReference" name="scalingRule"
        eType="#//ScalingRule" containment="true"/>
  </eClassifiers>

  <eClassifiers xsi:type="ecore:EClass" name="Server">
    <eStructuralFeatures xsi:type="ecore:EAttribute" name="name"
        eType="ecore:EDataType http://www.eclipse.org/emf/2002/Ecore#//EString"/>
    <eStructuralFeatures xsi:type="ecore:EAttribute" name="host"
        eType="ecore:EDataType http://www.eclipse.org/emf/2002/Ecore#//EString"/>
    <eStructuralFeatures xsi:type="ecore:EAttribute" name="port"
        eType="ecore:EDataType http://www.eclipse.org/emf/2002/Ecore#//EInt"/>
    <eStructuralFeatures xsi:type="ecore:EAttribute" name="weight"
        eType="ecore:EDataType http://www.eclipse.org/emf/2002/Ecore#//EInt"
        defaultValueLiteral="1"/>
    <eStructuralFeatures xsi:type="ecore:EAttribute" name="maxConnections"
        eType="ecore:EDataType http://www.eclipse.org/emf/2002/Ecore#//EInt"/>
    <eStructuralFeatures xsi:type="ecore:EAttribute" name="enabled"
        eType="ecore:EDataType http://www.eclipse.org/emf/2002/Ecore#//EBoolean"
        defaultValueLiteral="true"/>
  </eClassifiers>

  <eClassifiers xsi:type="ecore:EEnum" name="Algorithm">
    <eLiterals name="ROUND_ROBIN"/>
    <eLiterals name="WEIGHTED_ROUND_ROBIN" value="1"/>
    <eLiterals name="LEAST_CONNECTIONS" value="2"/>
    <eLiterals name="IP_HASH" value="3"/>
    <eLiterals name="RANDOM" value="4"/>
  </eClassifiers>

  <eClassifiers xsi:type="ecore:EEnum" name="Protocol">
    <eLiterals name="HTTP"/>
    <eLiterals name="HTTPS" value="1"/>
    <eLiterals name="TCP" value="2"/>
    <eLiterals name="UDP" value="3"/>
  </eClassifiers>

  <eClassifiers xsi:type="ecore:EEnum" name="Metric">
    <eLiterals name="CPU_USAGE"/>
    <eLiterals name="MEMORY_USAGE" value="1"/>
    <eLiterals name="REQUEST_COUNT" value="2"/>
    <eLiterals name="RESPONSE_TIME" value="3"/>
    <eLiterals name="ACTIVE_CONNECTIONS" value="4"/>
  </eClassifiers>

</ecore:EPackage>
\end{lstlisting}

% ============================================================
% CHAPITRE 6 : CONTRAINTES OCL
% ============================================================
\chapter{Contraintes OCL}

Chaque règle de gestion est formalisée en OCL (\textit{Object Constraint Language}) :

\section{RG01 -- Unicité des noms de serveurs}
\begin{lstlisting}[language=OCL]
context Cluster
inv uniqueServerNames:
  self.servers->isUnique(s | s.name)
\end{lstlisting}

\section{RG02 -- Port valide}
\begin{lstlisting}[language=OCL]
context Server
inv validPort:
  self.port >= 1 and self.port <= 65535

context Listener
inv validListenerPort:
  self.port >= 1 and self.port <= 65535
\end{lstlisting}

\section{RG03 -- Poids positif}
\begin{lstlisting}[language=OCL]
context Server
inv positiveWeight:
  self.weight >= 1
\end{lstlisting}

\section{RG04 -- Cluster non vide}
\begin{lstlisting}[language=OCL]
context Cluster
inv nonEmptyCluster:
  self.servers->notEmpty()
\end{lstlisting}

\section{RG05 -- Health check cohérent}
\begin{lstlisting}[language=OCL]
context HealthCheck
inv intervalGreaterThanTimeout:
  self.interval > self.timeout
\end{lstlisting}

\section{RG06 -- Seuils de scaling cohérents}
\begin{lstlisting}[language=OCL]
context ScalingRule
inv scaleUpGreaterThanDown:
  self.scaleUpThreshold > self.scaleDownThreshold
\end{lstlisting}

\section{RG07 -- Instances min/max cohérentes}
\begin{lstlisting}[language=OCL]
context ScalingRule
inv minLessOrEqualMax:
  self.minInstances <= self.maxInstances
\end{lstlisting}

\section{RG08 -- Listener unique par port}
\begin{lstlisting}[language=OCL]
context LoadBalancerSystem
inv uniqueListenerPorts:
  self.listeners->isUnique(l | l.port)
\end{lstlisting}

\section{RG09 -- Algorithme Weighted requiert des poids}
\begin{lstlisting}[language=OCL]
context Cluster
inv weightedRequiresWeights:
  self.loadBalancer.algorithm = Algorithm::WEIGHTED_ROUND_ROBIN
  implies self.servers->forAll(s | s.weight >= 1)
\end{lstlisting}

\section{RG10 -- Session Persistence cookie}
\begin{lstlisting}[language=OCL]
context SessionPersistence
inv cookieRequiresName:
  self.type = PersistenceType::COOKIE
  implies self.cookieName <> '' and self.cookieName <> null
\end{lstlisting}

\section{RG11 -- Au moins un listener}
\begin{lstlisting}[language=OCL]
context LoadBalancerSystem
inv atLeastOneListener:
  self.listeners->size() >= 1
\end{lstlisting}

\section{RG12 -- Serveur actif dans le cluster}
\begin{lstlisting}[language=OCL]
context Cluster
inv atLeastOneEnabledServer:
  self.servers->exists(s | s.enabled = true)
\end{lstlisting}

% ============================================================
% CHAPITRE 7 : NIVEAUX MDA
% ============================================================
\chapter{Niveaux MDA et architecture de l'approche}

\section{Les niveaux MDA pris en compte}

\begin{table}[H]
\centering
\caption{Niveaux MDA et artefacts produits}
\begin{tabular}{|p{2.5cm}|p{5cm}|p{5.5cm}|}
\hline
\textbf{Niveau} & \textbf{Description} & \textbf{Artefact produit} \\
\hline
CIM & Modèle indépendant de l'informatique & Description du domaine, concepts, règles \\
\hline
PIM & Modèle indépendant de la plateforme & Métamodèle Ecore, contraintes OCL, modèles DSL (.lb) \\
\hline
PSM & Modèle spécifique plateforme & Modèle intermédiaire Nginx/HAProxy (transfo ATL) \\
\hline
Code & Code exécutable généré & nginx.conf, docker-compose.yml (transfo Acceleo) \\
\hline
\end{tabular}
\end{table}

\section{Schéma d'architecture de l'approche}

\begin{figure}[H]
\centering
\begin{tikzpicture}[
  level/.style={rectangle, draw, thick, minimum width=14cm, minimum height=2cm, text centered},
  artifact/.style={rectangle, draw, fill=gray!10, minimum width=3cm, minimum height=0.8cm, text centered, font=\small},
  transfo/.style={->, >=stealth, thick, red!70!black},
  node distance=0.8cm
]

% CIM
\node[level, fill=blue!5] (cim) at (0,0) {};
\node[font=\bfseries] at (-5.5,0.6) {CIM};
\node[artifact] at (0,0) {Description domaine + Concepts + Règles};

% PIM
\node[level, fill=green!5] (pim) at (0,-3.5) {};
\node[font=\bfseries] at (-5.5,-2.5) {PIM};
\node[artifact] (ecore) at (-4,-3.5) {Ecore (.ecore)};
\node[artifact] (ocl) at (-0.5,-3.5) {OCL (.ocl)};
\node[artifact] (xtext) at (3,-3.5) {Xtext (.xtext)};
\node[artifact] (model) at (3,-4.5) {Modèle (.lb)};

% PSM
\node[level, fill=orange!5] (psm) at (0,-7) {};
\node[font=\bfseries] at (-5.5,-6.5) {PSM};
\node[artifact] (nginx_m) at (-2,-7) {Modèle Nginx};
\node[artifact] (haproxy_m) at (2.5,-7) {Modèle HAProxy};

% Code
\node[level, fill=red!5] (code) at (0,-10) {};
\node[font=\bfseries] at (-5.5,-9.5) {Code};
\node[artifact] (nginxconf) at (-4,-10) {nginx.conf};
\node[artifact] (docker) at (0,-10) {docker-compose.yml};
\node[artifact] (haproxyconf) at (4,-10) {haproxy.cfg};

% Transformations
\draw[transfo] (0,-1) -- node[right, font=\small\itshape] {Analyse} (0,-2.2);
\draw[transfo] (model.south) -- node[right, font=\small\itshape] {ATL (M2M)} (0,-5.8);
\draw[transfo] (nginx_m.south) -- node[left, font=\small\itshape] {Acceleo (M2T)} (-2,-8.8);
\draw[transfo] (haproxy_m.south) -- node[right, font=\small\itshape] {Acceleo (M2T)} (2.5,-8.8);

\end{tikzpicture}
\caption{Architecture MDA de l'approche}
\end{figure}

\section{Justification des choix technologiques}

\begin{table}[H]
\centering
\caption{Outils et technologies utilisés}
\begin{tabular}{|p{3cm}|p{4cm}|p{6cm}|}
\hline
\textbf{Outil} & \textbf{Rôle} & \textbf{Justification} \\
\hline
EMF / Ecore & Métamodèle (syntaxe abstraite) & Standard pour la méta-modélisation Eclipse \\
\hline
OCL & Contraintes & Langage standardisé OMG pour contraintes \\
\hline
Xtext & Syntaxe concrète DSL & Framework mature, éditeur intégré Eclipse \\
\hline
ATL & Transformations M2M & Langage hybride déclaratif/impératif \\
\hline
Acceleo & Transformations M2T & Génération basée templates, compatible EMF \\
\hline
\end{tabular}
\end{table}

% ============================================================
% CHAPITRE 8 : TRANSFORMATIONS
% ============================================================
\chapter{Spécification des transformations}

\section{Vue d'ensemble}

\begin{table}[H]
\centering
\caption{Transformations définies}
\begin{tabular}{|p{2.5cm}|p{2cm}|p{3cm}|p{3cm}|p{1.5cm}|}
\hline
\textbf{Transfo} & \textbf{Type} & \textbf{Source} & \textbf{Cible} & \textbf{Outil} \\
\hline
T1: PIM$\rightarrow$PSM Nginx & M2M & Modèle .lb & Modèle Nginx & ATL \\
\hline
T2: PIM$\rightarrow$PSM HAProxy & M2M & Modèle .lb & Modèle HAProxy & ATL \\
\hline
T3: PSM Nginx$\rightarrow$Code & M2T & Modèle Nginx & nginx.conf & Acceleo \\
\hline
T4: PSM HAProxy$\rightarrow$Code & M2T & Modèle HAProxy & haproxy.cfg & Acceleo \\
\hline
T5: PIM$\rightarrow$Docker & M2T & Modèle .lb & docker-compose.yml & Acceleo \\
\hline
\end{tabular}
\end{table}

\section{Définition de la grammaire Xtext (syntaxe concrète)}

La grammaire Xtext définit la syntaxe concrète textuelle de notre DSL. Elle permet de créer un éditeur avec coloration syntaxique, auto-complétion et validation intégrée dans Eclipse. Le fichier de grammaire est \texttt{LoadBalancer.xtext} :

\begin{lstlisting}[language=Xtext, caption={Grammaire Xtext complète du DSL de répartition de charge}]
grammar org.loadbalancer.dsl.LoadBalancer
  with org.eclipse.xtext.common.Terminals

generate loadBalancer "http://www.loadbalancer.org/dsl/LoadBalancer"

// ==========================================
// Racine du modele
// ==========================================
LoadBalancerSystem:
  'LoadBalancerSystem' name=STRING '{'
    ('description:' description=STRING)?
    ('version:' version=STRING)?
    clusters+=Cluster+
    listeners+=Listener+
    (alerts+=Alert)*
  '}';

// ==========================================
// Cluster de serveurs
// ==========================================
Cluster:
  'cluster' name=ID '{'
    loadBalancer=LoadBalancerConfig
    servers+=Server+
    (healthCheck=HealthCheck)?
    (scalingRule=ScalingRule)?
  '}';

// ==========================================
// Configuration du repartiteur
// ==========================================
LoadBalancerConfig:
  'loadbalancer' '{'
    'algorithm:' algorithm=Algorithm
    'stickySession:' stickySession=BooleanValue
    (sessionPersistence=SessionPersistence)?
  '}';

// ==========================================
// Serveur
// ==========================================
Server:
  'server' name=ID '{'
    'host:' host=STRING
    'port:' port=INT
    ('weight:' weight=INT)?
    ('maxConnections:' maxConnections=INT)?
    'enabled:' enabled=BooleanValue
  '}';

// ==========================================
// Health Check
// ==========================================
HealthCheck:
  'healthCheck' '{'
    'protocol:' protocol=Protocol
    'path:' path=STRING
    'interval:' interval=INT
    'timeout:' timeout=INT
    ('healthyThreshold:' healthyThreshold=INT)?
    ('unhealthyThreshold:' unhealthyThreshold=INT)?
  '}';

// ==========================================
// Persistence de session
// ==========================================
SessionPersistence:
  'sessionPersistence' '{'
    'type:' type=PersistenceType
    ('cookieName:' cookieName=STRING)?
    ('ttl:' ttl=INT)?
  '}';

// ==========================================
// Regle de mise a l'echelle
// ==========================================
ScalingRule:
  'scalingRule' name=ID '{'
    'metric:' metric=Metric
    'scaleUpThreshold:' scaleUpThreshold=DOUBLE
    'scaleDownThreshold:' scaleDownThreshold=DOUBLE
    'minInstances:' minInstances=INT
    'maxInstances:' maxInstances=INT
  '}';

// ==========================================
// Listener (point d'entree)
// ==========================================
Listener:
  'listener' name=ID '{'
    'protocol:' protocol=Protocol
    'port:' port=INT
    'targetCluster:' targetCluster=[Cluster]
  '}';

// ==========================================
// Alerte
// ==========================================
Alert:
  'alert' name=ID '{'
    'metric:' metric=Metric
    'threshold:' threshold=DOUBLE
    'action:' action=STRING
  '}';

// ==========================================
// Enumerations
// ==========================================
enum Algorithm:
  ROUND_ROBIN = 'ROUND_ROBIN' |
  WEIGHTED_ROUND_ROBIN = 'WEIGHTED_ROUND_ROBIN' |
  LEAST_CONNECTIONS = 'LEAST_CONNECTIONS' |
  IP_HASH = 'IP_HASH' |
  RANDOM = 'RANDOM';

enum Protocol:
  HTTP = 'HTTP' |
  HTTPS = 'HTTPS' |
  TCP = 'TCP';

enum Metric:
  CPU_USAGE = 'CPU_USAGE' |
  MEMORY_USAGE = 'MEMORY_USAGE' |
  RESPONSE_TIME = 'RESPONSE_TIME' |
  REQUEST_COUNT = 'REQUEST_COUNT' |
  ERROR_RATE = 'ERROR_RATE';

enum PersistenceType:
  COOKIE = 'COOKIE' |
  IP = 'IP' |
  HEADER = 'HEADER';

// ==========================================
// Types terminaux
// ==========================================
BooleanValue:
  value=('true' | 'false');

terminal DOUBLE returns ecore::EDouble:
  INT '.' INT;
\end{lstlisting}

\subsection{Explication de la grammaire}

La grammaire ci-dessus définit la syntaxe concrète complète du DSL. Voici les éléments clés :

\begin{itemize}
  \item \textbf{LoadBalancerSystem} : élément racine, contient les clusters, listeners et alertes.
  \item \textbf{Cluster} : regroupe un ensemble de serveurs avec une stratégie de répartition commune. Chaque cluster possède une configuration de \texttt{LoadBalancerConfig}, au moins un \texttt{Server}, et optionnellement un \texttt{HealthCheck} et un \texttt{ScalingRule}.
  \item \textbf{Server} : représente un serveur backend avec son hôte, port, poids et état d'activation.
  \item \textbf{Listener} : définit un point d'écoute (port d'entrée) et le cluster cible via une \textbf{référence croisée} (\texttt{[Cluster]}).
  \item \textbf{Énumérations} : \texttt{Algorithm}, \texttt{Protocol}, \texttt{Metric} et \texttt{PersistenceType} contraignent les valeurs acceptées.
  \item \textbf{Terminal DOUBLE} : terminal personnalisé pour les valeurs décimales (seuils de scaling et alertes).
\end{itemize}

\subsection{Artefacts générés par Xtext}

L'exécution de \texttt{Generate Xtext Artifacts} produit automatiquement :

\begin{table}[H]
\centering
\caption{Artefacts générés par le workflow Xtext}
\begin{tabular}{|p{5cm}|p{8cm}|}
\hline
\textbf{Artefact} & \textbf{Description} \\
\hline
Parseur ANTLR & Analyse syntaxique des fichiers \texttt{.lb} \\
\hline
Modèle EMF (EPackage) & Syntaxe abstraite dérivée de la grammaire \\
\hline
Éditeur Eclipse & Coloration syntaxique, auto-complétion, outline \\
\hline
Sérialiseur & Conversion modèle $\leftrightarrow$ texte \\
\hline
Validateur & Classe \texttt{LoadBalancerValidator.xtend} pour les contraintes \\
\hline
Formateur & Indentation et formatage du code DSL \\
\hline
Scope Provider & Résolution des références croisées (ex: \texttt{targetCluster}) \\
\hline
\end{tabular}
\end{table}

\section{Transformation T1 : PIM $\rightarrow$ PSM Nginx (ATL)}

\begin{lstlisting}[language=ATL, caption={Transformation ATL PIM vers Nginx}]
-- Fichier : PIM2Nginx.atl
module PIM2Nginx;
create OUT : NginxModel from IN : LoadBalancerModel;

rule System2NginxConfig {
  from s : LoadBalancerModel!LoadBalancerSystem
  to   c : NginxModel!NginxConfig (
    workerProcesses <- 'auto',
    workerConnections <- 1024,
    upstreams <- s.clusters->collect(cl |
      thisModule.Cluster2Upstream(cl)),
    servers <- s.listeners->collect(l |
      thisModule.Listener2ServerBlock(l))
  )
}

lazy rule Cluster2Upstream {
  from cl : LoadBalancerModel!Cluster
  to   u : NginxModel!Upstream (
    name <- cl.name,
    algorithm <- cl.loadBalancer.algorithm.mapNginxAlgo(),
    servers <- cl.servers->select(s | s.enabled)
              ->collect(s | thisModule.Server2UpstreamServer(s))
  )
}

lazy rule Server2UpstreamServer {
  from s : LoadBalancerModel!Server
  to   us : NginxModel!UpstreamServer (
    address <- s.host + ':' + s.port.toString(),
    weight <- s.weight
  )
}

lazy rule Listener2ServerBlock {
  from l : LoadBalancerModel!Listener
  to   sb : NginxModel!ServerBlock (
    listenPort <- l.port,
    proxyPass <- 'http://' + l.targetCluster.name
  )
}

helper context LoadBalancerModel!Algorithm
def : mapNginxAlgo() : String =
  if self = #LEAST_CONNECTIONS then 'least_conn'
  else if self = #IP_HASH then 'ip_hash'
  else '' endif endif;
\end{lstlisting}

\section{Transformation T3 : PSM Nginx $\rightarrow$ Code (Acceleo)}

\begin{lstlisting}[caption={Template Acceleo pour nginx.conf}]
[comment encoding = UTF-8 /]
[module generateNginx('http://www.example.org/nginx')]

[template public generateNginxConf(config : NginxConfig)]
[comment @main /]
[file ('nginx.conf', false, 'UTF-8')]
worker_processes [config.workerProcesses/];

events {
    worker_connections [config.workerConnections/];
}

http {
[for (upstream : Upstream | config.upstreams)]
    upstream [upstream.name/] {
[if (upstream.algorithm = 'least_conn')]
        least_conn;
[elseif (upstream.algorithm = 'ip_hash')]
        ip_hash;
[/if]
[for (server : UpstreamServer | upstream.servers)]
        server [server.address/] weight=[server.weight/];
[/for]
    }

[/for]
[for (server : ServerBlock | config.servers)]
    server {
        listen [server.listenPort/];
        location / {
            proxy_pass [server.proxyPass/];
        }
    }
[/for]
}
[/file]
[/template]
\end{lstlisting}

\section{Transformation T5 : PIM $\rightarrow$ Docker Compose (Acceleo)}

\begin{lstlisting}[caption={Template Acceleo pour docker-compose.yml}]
[comment encoding = UTF-8 /]
[module generateDockerCompose('http://www.example.org/loadbalancer')]

[template public generateCompose(system : LoadBalancerSystem)]
[comment @main /]
[file ('docker-compose.yml', false, 'UTF-8')]
version: '3.8'
services:
[for (cluster : Cluster | system.clusters)]
[for (server : Server | cluster.servers)]
  [server.name/]:
    image: nginx:alpine
    ports:
      - "[server.port/]:[server.port/]"
[/for]
  [cluster.name/]-lb:
    image: nginx:alpine
    volumes:
      - ./nginx.conf:/etc/nginx/nginx.conf:ro
    ports:
[for (listener : Listener | system.listeners)]
[if (listener.targetCluster.name = cluster.name)]
      - "[listener.port/]:[listener.port/]"
[/if]
[/for]
    depends_on:
[for (server : Server | cluster.servers)]
      - [server.name/]
[/for]
[/for]
[/file]
[/template]
\end{lstlisting}

% ============================================================
% CHAPITRE 9 : INSTALLATION DES OUTILS
% ============================================================
\chapter{Installation et configuration des outils}

Cette section décrit l'installation des outils nécessaires pour exécuter et tester les artefacts générés par notre chaîne de transformation (nginx.conf et docker-compose.yml).

\section{Prérequis système}

\begin{itemize}
  \item Système d'exploitation : Linux (Ubuntu/Debian recommandé), macOS ou Windows avec WSL2
  \item Minimum 4 Go de RAM disponible
  \item Connexion internet pour télécharger les images Docker
\end{itemize}

\section{Installation de Docker}

Docker est le moteur de conteneurisation qui permet d'exécuter les serveurs backend et le load balancer.

\subsection{Sur Ubuntu/Debian}

\begin{lstlisting}[language=bash, caption={Installation de Docker sur Ubuntu/Debian}]
# 1. Mise a jour du systeme
sudo apt update && sudo apt upgrade -y

# 2. Installation des dependances
sudo apt install -y apt-transport-https ca-certificates \
  curl gnupg lsb-release

# 3. Ajout de la cle GPG officielle Docker
curl -fsSL https://download.docker.com/linux/ubuntu/gpg | \
  sudo gpg --dearmor -o /usr/share/keyrings/docker-archive-keyring.gpg

# 4. Ajout du depot Docker
echo "deb [arch=$(dpkg --print-architecture) \
  signed-by=/usr/share/keyrings/docker-archive-keyring.gpg] \
  https://download.docker.com/linux/ubuntu \
  $(lsb_release -cs) stable" | \
  sudo tee /etc/apt/sources.list.d/docker.list > /dev/null

# 5. Installation de Docker Engine
sudo apt update
sudo apt install -y docker-ce docker-ce-cli containerd.io

# 6. Ajout de l'utilisateur au groupe docker (evite sudo)
sudo usermod -aG docker $USER
newgrp docker

# 7. Verification
docker --version
docker run hello-world
\end{lstlisting}

\subsection{Sur macOS}

\begin{lstlisting}[language=bash, caption={Installation de Docker sur macOS}]
# Telecharger Docker Desktop depuis :
# https://www.docker.com/products/docker-desktop/
# Ou via Homebrew :
brew install --cask docker

# Lancer Docker Desktop puis verifier :
docker --version
\end{lstlisting}

\section{Installation de Docker Compose}

Docker Compose V2 est inclus avec Docker Desktop. Sur Linux, il peut être installé séparément :

\begin{lstlisting}[language=bash, caption={Installation de Docker Compose sur Linux}]
# Docker Compose V2 (plugin Docker)
sudo apt install -y docker-compose-plugin

# Verification
docker compose version

# Alternative : installation standalone
sudo curl -SL https://github.com/docker/compose/releases/latest/download/docker-compose-linux-x86_64 \
  -o /usr/local/bin/docker-compose
sudo chmod +x /usr/local/bin/docker-compose
docker-compose --version
\end{lstlisting}

\section{Installation de Nginx (pour tests locaux sans Docker)}

Bien que notre approche utilise Docker, il est utile d'avoir Nginx installé localement pour des tests rapides :

\begin{lstlisting}[language=bash, caption={Installation de Nginx}]
# Ubuntu/Debian
sudo apt install -y nginx

# Verification
nginx -v
sudo systemctl status nginx

# Test de configuration
sudo nginx -t

# Demarrage
sudo systemctl start nginx
sudo systemctl enable nginx
\end{lstlisting}

\section{Installation des outils IDM (Eclipse)}

Pour la partie modélisation et transformation :

\begin{lstlisting}[language=bash, caption={Installation d'Eclipse Modeling Tools}]
# 1. Telecharger Eclipse Modeling Tools depuis :
#    https://www.eclipse.org/downloads/packages/
#    Choisir "Eclipse Modeling Tools"

# 2. Installer les plugins via Help > Eclipse Marketplace :
#    - Xtext Complete SDK
#    - ATL/EMFTVM
#    - Acceleo
#    - OCL Tools (Eclipse OCL)

# 3. Ou via Help > Install New Software :
#    URL: https://download.eclipse.org/releases/2023-12
#    Selectionner : Modeling > EMF, Xtext, ATL, Acceleo, OCL
\end{lstlisting}

\section{Outils complémentaires utiles}

\begin{lstlisting}[language=bash, caption={Outils complémentaires}]
# curl - pour tester les requetes HTTP
sudo apt install -y curl

# jq - pour formatter les reponses JSON
sudo apt install -y jq

# htop - pour surveiller les ressources
sudo apt install -y htop

# ab (Apache Bench) - pour les tests de charge
sudo apt install -y apache2-utils

# wrk - outil de benchmark HTTP performant
sudo apt install -y wrk
\end{lstlisting}

\section{Vérification de l'installation complète}

\begin{lstlisting}[language=bash, caption={Script de vérification}]
#!/bin/bash
echo "=== Verification de l'installation ==="
echo ""
echo -n "Docker: "
docker --version 2>/dev/null || echo "NON INSTALLE"
echo -n "Docker Compose: "
docker compose version 2>/dev/null || echo "NON INSTALLE"
echo -n "Nginx: "
nginx -v 2>&1 || echo "NON INSTALLE"
echo -n "curl: "
curl --version | head -1 || echo "NON INSTALLE"
echo ""
echo "=== Test Docker ==="
docker run --rm hello-world 2>/dev/null && echo "Docker OK" || echo "Docker ERREUR"
\end{lstlisting}

% ============================================================
% CHAPITRE 10 : ETUDE DE CAS
% ============================================================
\chapter{Étude de cas : Application E-Commerce}

\section{Description du cas d'étude}

Nous illustrons notre approche avec un cas concret : une application web e-commerce déployée sur un système distribué. L'application possède :

\begin{itemize}
  \item Un cluster de 3 serveurs web (frontend) avec répartition Round Robin
  \item Un cluster de 2 serveurs API (backend) avec répartition Least Connections
  \item Des health checks HTTP sur chaque cluster
  \item Une règle de scaling automatique basée sur l'utilisation CPU
  \item Des alertes en cas de dépassement de seuils
\end{itemize}

\section{Modèle DSL du cas d'étude}

Voici le modèle écrit avec notre DSL (syntaxe concrète textuelle) :

\begin{lstlisting}[language=DSL, caption={Modèle DSL de l'application e-commerce (ecommerce.lb)}]
LoadBalancerSystem "ECommerceSystem" {
  description: "Systeme de repartition de charge pour e-commerce"
  version: "1.0"

  cluster WebFrontend {
    loadbalancer {
      algorithm: ROUND_ROBIN
      stickySession: true
      sessionPersistence {
        type: COOKIE
        cookieName: "SESSIONID"
        ttl: 3600
      }
    }

    server web1 {
      host: "192.168.1.10"
      port: 8080
      weight: 3
      maxConnections: 1000
      enabled: true
    }
    server web2 {
      host: "192.168.1.11"
      port: 8080
      weight: 2
      maxConnections: 800
      enabled: true
    }
    server web3 {
      host: "192.168.1.12"
      port: 8080
      weight: 1
      maxConnections: 500
      enabled: true
    }

    healthCheck {
      protocol: HTTP
      path: "/health"
      interval: 30
      timeout: 10
      healthyThreshold: 3
      unhealthyThreshold: 2
    }

    scalingRule autoScaleWeb {
      metric: CPU_USAGE
      scaleUpThreshold: 80.0
      scaleDownThreshold: 30.0
      minInstances: 2
      maxInstances: 10
    }
  }

  cluster APIBackend {
    loadbalancer {
      algorithm: LEAST_CONNECTIONS
      stickySession: false
    }

    server api1 {
      host: "192.168.2.10"
      port: 3000
      weight: 1
      maxConnections: 500
      enabled: true
    }
    server api2 {
      host: "192.168.2.11"
      port: 3000
      weight: 1
      maxConnections: 500
      enabled: true
    }

    healthCheck {
      protocol: HTTP
      path: "/api/health"
      interval: 15
      timeout: 5
      healthyThreshold: 2
      unhealthyThreshold: 3
    }
  }

  listener httpListener {
    protocol: HTTP
    port: 80
    targetCluster: WebFrontend
  }

  listener apiListener {
    protocol: HTTP
    port: 3000
    targetCluster: APIBackend
  }

  alert highCPU {
    metric: CPU_USAGE
    threshold: 90.0
    action: "email:admin@ecommerce.com"
  }

  alert highLatency {
    metric: RESPONSE_TIME
    threshold: 5000.0
    action: "webhook:https://hooks.slack.com/alert"
  }
}
\end{lstlisting}

\section{Code généré : nginx.conf}

Après exécution de la chaîne de transformations, le fichier nginx.conf suivant est généré :

\begin{lstlisting}[language={}, caption={Fichier nginx.conf généré}]
worker_processes auto;

events {
    worker_connections 1024;
}

http {
    upstream WebFrontend {
        server 192.168.1.10:8080 weight=3;
        server 192.168.1.11:8080 weight=2;
        server 192.168.1.12:8080 weight=1;
    }

    upstream APIBackend {
        least_conn;
        server 192.168.2.10:3000 weight=1;
        server 192.168.2.11:3000 weight=1;
    }

    server {
        listen 80;
        location / {
            proxy_pass http://WebFrontend;
        }
    }

    server {
        listen 3000;
        location / {
            proxy_pass http://APIBackend;
        }
    }
}
\end{lstlisting}

\section{Code généré : docker-compose.yml}

\begin{lstlisting}[language={}, caption={Fichier docker-compose.yml généré}]
version: '3.8'

services:
  web1:
    build: ./servers/web
    environment:
      - SERVER_NAME=web1
      - SERVER_PORT=8080
    healthcheck:
      test: ["CMD", "curl", "-f", "http://localhost:8080/health"]
      interval: 30s
      timeout: 10s
      retries: 3

  web2:
    build: ./servers/web
    environment:
      - SERVER_NAME=web2
      - SERVER_PORT=8080
    healthcheck:
      test: ["CMD", "curl", "-f", "http://localhost:8080/health"]
      interval: 30s
      timeout: 10s
      retries: 3

  web3:
    build: ./servers/web
    environment:
      - SERVER_NAME=web3
      - SERVER_PORT=8080
    healthcheck:
      test: ["CMD", "curl", "-f", "http://localhost:8080/health"]
      interval: 30s
      timeout: 10s
      retries: 3

  api1:
    build: ./servers/api
    environment:
      - SERVER_NAME=api1
      - SERVER_PORT=3000
    healthcheck:
      test: ["CMD", "curl", "-f", "http://localhost:3000/api/health"]
      interval: 15s
      timeout: 5s
      retries: 3

  api2:
    build: ./servers/api
    environment:
      - SERVER_NAME=api2
      - SERVER_PORT=3000
    healthcheck:
      test: ["CMD", "curl", "-f", "http://localhost:3000/api/health"]
      interval: 15s
      timeout: 5s
      retries: 3

  loadbalancer:
    image: nginx:alpine
    volumes:
      - ./nginx.conf:/etc/nginx/nginx.conf:ro
    ports:
      - "80:80"
      - "3000:3000"
    depends_on:
      - web1
      - web2
      - web3
      - api1
      - api2
\end{lstlisting}

\section{Déploiement et test}

\subsection{Lancement de l'infrastructure}

\begin{lstlisting}[language=bash, caption={Commandes de déploiement}]
# Se placer dans le dossier du cas d'etude
cd case-study/

# Construire et demarrer tous les services
docker compose up --build -d

# Verifier que tous les conteneurs sont actifs
docker compose ps

# Voir les logs
docker compose logs -f
\end{lstlisting}

\subsection{Tests de la répartition de charge}

\begin{lstlisting}[language=bash, caption={Tests de vérification}]
# Test du cluster Web Frontend (Round Robin)
for i in $(seq 1 10); do
  echo "Request $i:"
  curl -s http://localhost:80/ | jq .
  echo "---"
done

# Test du cluster API Backend (Least Connections)
for i in $(seq 1 10); do
  echo "Request $i:"
  curl -s http://localhost:3000/api/health | jq .
  echo "---"
done

# Test de charge avec wrk
wrk -t4 -c100 -d30s http://localhost:80/

# Test de charge avec ab
ab -n 1000 -c 50 http://localhost:80/
\end{lstlisting}

\subsection{Vérification du health check}

\begin{lstlisting}[language=bash, caption={Test du health check}]
# Arreter un serveur pour tester le failover
docker compose stop web3

# Verifier que les requetes sont redistribuees
for i in $(seq 1 6); do
  curl -s http://localhost:80/ | jq .serverName
done

# Redemarrer le serveur
docker compose start web3
\end{lstlisting}

% ============================================================
% CHAPITRE 11 : TUTORIEL
% ============================================================
\chapter{Tutoriel de réalisation pas à pas}

\section{Prérequis}

\begin{itemize}
  \item Eclipse Modeling Tools (version 2023-12 ou supérieure)
  \item Plugins : EMF, Xtext, ATL, Acceleo, OCL Tools
  \item Java JDK 17+
  \item Docker et Docker Compose (pour les tests)
\end{itemize}

\section{Étape 1 : Création du métamodèle Ecore}

\begin{enumerate}
  \item Dans Eclipse : \texttt{File $\rightarrow$ New $\rightarrow$ Other $\rightarrow$ Eclipse Modeling Framework $\rightarrow$ Ecore Modeling Project}
  \item Nommer le projet : \texttt{org.loadbalancer.metamodel}
  \item Créer le fichier \texttt{loadbalancer.ecore}
  \item Définir toutes les classes, attributs et relations (voir chapitre 5)
  \item Ouvrir avec l'éditeur graphique pour visualiser le diagramme
\end{enumerate}

\section{Étape 2 : Génération du code Java EMF}

\begin{enumerate}
  \item Clic droit sur \texttt{loadbalancer.ecore} $\rightarrow$ \texttt{New} $\rightarrow$ \texttt{EMF Generator Model}
  \item Créer \texttt{loadbalancer.genmodel}
  \item Ouvrir le \texttt{.genmodel} $\rightarrow$ Clic droit sur la racine $\rightarrow$ \texttt{Generate All}
\end{enumerate}

\section{Étape 3 : Définition de la syntaxe concrète (Xtext)}

\begin{enumerate}
  \item \texttt{File $\rightarrow$ New $\rightarrow$ Project $\rightarrow$ Xtext $\rightarrow$ Xtext Project}
  \item Nommer : \texttt{org.loadbalancer.dsl}, extension : \texttt{.lb}
  \item Définir la grammaire (voir chapitre 8, section Xtext)
  \item Exécuter : Clic droit $\rightarrow$ \texttt{Run As $\rightarrow$ Generate Xtext Artifacts}
  \item Tester dans une Eclipse Application
\end{enumerate}

\section{Étape 4 : Implémentation des contraintes}

\begin{enumerate}
  \item Ouvrir le fichier \texttt{LoadBalancerValidator.xtend} généré par Xtext
  \item Ajouter les méthodes de validation (voir chapitre 6)
  \item Chaque méthode annotée \texttt{@Check} implémente une règle OCL
\end{enumerate}

\section{Étape 5 : Transformations ATL}

\begin{enumerate}
  \item \texttt{File $\rightarrow$ New $\rightarrow$ ATL Project} : \texttt{org.loadbalancer.transformations}
  \item Créer le métamodèle cible \texttt{nginx.ecore}
  \item Écrire \texttt{PIM2Nginx.atl} (voir chapitre 8)
  \item Configurer le \texttt{Run Configuration ATL}
\end{enumerate}

\section{Étape 6 : Génération de code (Acceleo)}

\begin{enumerate}
  \item \texttt{File $\rightarrow$ New $\rightarrow$ Acceleo Project} : \texttt{org.loadbalancer.generators}
  \item Créer les templates \texttt{.mtl} (voir chapitre 8)
  \item Configurer le launch configuration
\end{enumerate}

\section{Étape 7 : Test de bout en bout}

\begin{enumerate}
  \item Lancer une Eclipse Application
  \item Créer un fichier \texttt{ecommerce.lb}
  \item Vérifier la validation des contraintes
  \item Exécuter les transformations
  \item Récupérer les fichiers générés
  \item Tester avec Docker Compose (chapitre 10)
\end{enumerate}

% ============================================================
% CHAPITRE 12 : PERSPECTIVES
% ============================================================
\chapter{Perspectives et améliorations}

Ce qui resterait à faire pour atteindre pleinement les objectifs :

\begin{itemize}
  \item Support complet HAProxy comme cible PSM supplémentaire
  \item Génération de manifestes Kubernetes (Deployment, Service, Ingress)
  \item Éditeur graphique avec Eclipse Sirius
  \item Support du load balancing DNS
  \item Intégration monitoring (Prometheus + Grafana)
  \item Validation dynamique avec simulation de charge
  \item Support multi-cloud (AWS ALB, GCP LB, Azure LB)
  \item Interface web pour les non-développeurs
\end{itemize}

% ============================================================
% CONCLUSION
% ============================================================
\chapter*{Conclusion}
\addcontentsline{toc}{chapter}{Conclusion}

Ce projet a permis de mettre en pratique les concepts fondamentaux de l'Ingénierie Dirigée par les Modèles (IDM) à travers la conception d'un DSL dédié à la répartition de charge dans les systèmes distribués.

Nous avons suivi une démarche MDA complète :
\begin{itemize}
  \item \textbf{Niveau CIM} : analyse et description du domaine de la répartition de charge
  \item \textbf{Niveau PIM} : conception du métamodèle Ecore, définition des contraintes OCL, et création de la syntaxe concrète Xtext
  \item \textbf{Niveau PSM} : transformation ATL vers des modèles spécifiques (Nginx)
  \item \textbf{Niveau Code} : génération automatique de fichiers de configuration via Acceleo
\end{itemize}

L'approche proposée démontre les avantages de l'IDM : abstraction, réutilisabilité, portabilité entre plateformes, et validation formelle des modèles. Le DSL créé permet aux administrateurs système de décrire leur infrastructure de répartition de charge de manière déclarative et indépendante de la technologie, tout en bénéficiant d'une génération automatique de code fiable et cohérente.

L'étude de cas avec Docker Compose et Nginx prouve la faisabilité de l'approche en permettant un déploiement réel de l'infrastructure générée.

% ============================================================
% REFERENCES
% ============================================================
\chapter*{Références bibliographiques}
\addcontentsline{toc}{chapter}{Références bibliographiques}

\begin{enumerate}[label={[\arabic*]}]
  \item J. Bézivin, \og On the Unification Power of Models \fg{}, \textit{Software and Systems Modeling}, vol.~4, no.~2, pp.~171--188, 2005.
  \item OMG, \og MDA Guide Version 2.0 \fg{}, Object Management Group, 2014.
  \item F. Jouault, F. Allilaire, J. Bézivin, I. Kurtev, \og ATL: A model transformation tool \fg{}, \textit{Science of Computer Programming}, vol.~72, no.~1-2, pp.~31--39, 2008.
  \item M. Bettini, \og Implementing Domain-Specific Languages with Xtext and Xtend \fg{}, Packt Publishing, 2\textsuperscript{nd} edition, 2016.
  \item Eclipse Foundation, \og Acceleo User Guide \fg{}, \url{https://wiki.eclipse.org/Acceleo}.
  \item Eclipse Foundation, \og Eclipse Modeling Framework (EMF) \fg{}, \url{https://www.eclipse.org/modeling/emf/}.
  \item Nginx Documentation, \og HTTP Load Balancing \fg{}, \url{https://docs.nginx.com/}.
  \item HAProxy Documentation, \og Configuration Manual \fg{}, \url{https://www.haproxy.org/}.
  \item T. Stahl, M. Voelter, \og Model-Driven Software Development \fg{}, Wiley, 2006.
  \item A. Kleppe, J. Warmer, W. Bast, \og MDA Explained \fg{}, Addison-Wesley, 2003.
\end{enumerate}

\end{document}
